\section{Limitations and Conclusion}
This study establishes a reproducible physical contact-RL basis and a limited historical residual-control geometry result. It does not establish a novel outcome-aware algorithm: the new effect-history pilot fails its promotion criterion against matched task-frame residual PPO. Each pilot model has only four rollout updates and 16 completed training episodes; the bounded null result does not disprove execution-history RL in general. The historical three-seed geometry result has substantial variance and limited P2 improvement. The new near-ceiling development set does not establish unseen-condition generalization, while direct PPO needs further matched tuning and rotation/shared-training ablations remain untested. Historical aggregate improvement must not be represented as a solved narrow-slot task.

State-based simulator observations include privileged pose and known dynamics context. A deployed perception system, visual/point-cloud actor, measured tool dynamics, true grasping, full robot CCD and real-robot validation have not been demonstrated. The synthetic cube and research paddle are not the field tray. Field references are kept private and separate. Long-horizon option interfaces, VLA instructions and skill recovery are future extensions, not completed contributions.

The next decision is deliberately bounded: do not add history channels or extend training on the near-ceiling pilot. First establish physical authority of a re-contact/contact-location parameter under cold-start paired executions, then compare the same strong feedback and residual controls in physically justified, non-saturated contact conditions. Promote only a repeatable P2 gain with a valid mechanism ablation. Negative results narrow this mechanism without discarding the platform or the requirement to study RL. A publishable application-method paper still requires an independently supported increment, more seeds, sealed evaluation groups and stronger related-method comparisons.
